% =========================================================
%  metadata.tex
%  Metadata buku dan identitas mata kuliah.
%  Diubah di sini akan berlaku pada seluruh dokumen.
% =========================================================

% --- Buku ------------------------------------------------
\newcommand{\BookTitle}{Pengantar Pemrograman Game}
\newcommand{\BookSubtitle}{Buku Ajar Berbasis \istilah{Outcome-Based Education} (OBE)}
\newcommand{\BookEdition}{Edisi Pertama}
\newcommand{\BookYear}{2026}
\newcommand{\BookAuthor}{Cecep Suwanda, S.Si., M.Kom.}

% --- Identitas mata kuliah -------------------------------
\newcommand{\NamaMK}{Pengantar Pemrograman Game}
\newcommand{\KodeMK}{MWCA009}
\newcommand{\BobotSKS}{3 SKS (2 SKS teori, 1 SKS praktikum)}
\newcommand{\SemesterMK}{Ganjil}
\newcommand{\TahunAkademik}{2026/2027}
\newcommand{\MataKuliahPrasyarat}{Algoritma dan Pemrograman}

% --- Identitas institusi ---------------------------------
\newcommand{\NamaProdi}{Teknik Informatika}
\newcommand{\NamaFakultas}{Teknologi Informasi}
\newcommand{\NamaPT}{Universitas Bale Bandung}
\newcommand{\NamaDosen}{Cecep Suwanda, S.Si., M.Kom.}

% --- Metadata berkas PDF ---------------------------------
\hypersetup{
  pdftitle  = {\BookTitle\ --- Buku Ajar \NamaMK},
  pdfauthor = {\BookAuthor},
  pdfsubject = {Buku ajar berbasis OBE untuk mata kuliah \NamaMK\ (\KodeMK)},
  pdfkeywords = {pemrograman game, Pygame, Unity, OBE, \NamaProdi}
}
